\documentclass[12pt]{article}
\usepackage[a4paper,tmargin=3cm,lmargin=2.5cm,rmargin=2.5cm,bmargin=3cm]{geometry}
\usepackage[utf8]{inputenc}
\usepackage[T1]{fontenc}
\usepackage[brazilian]{babel}
\usepackage[authoryear,round]{natbib}
\let\cite\citep
\let\citeonline\citet
\usepackage{amsmath,amssymb}
\usepackage{graphicx}
\usepackage{booktabs,array,longtable}
\usepackage{tikz}
\usetikzlibrary{arrows.meta,positioning}
\usepackage{fancyhdr}
\usepackage{xurl}
\usepackage{hyperref}
\hypersetup{colorlinks=true,linkcolor=black,filecolor=magenta,urlcolor=blue,citecolor=black,
pdftitle={EduAgent-OS: Fase 2 - Metodologia},
pdfauthor={Rafael Attilio Agricola; Jociclelio Castro Macedo Junior}}
\setlength{\emergencystretch}{3em}
\setlength{\headheight}{14.5pt}
\setlength{\footskip}{32pt}
\pagestyle{fancy}
\fancyhf{}
\renewcommand{\headrulewidth}{0pt}
\fancyfoot[C]{IC/UNICAMP \\ 2026}
\newcommand{\code}[1]{\nolinkurl{#1}}
\newcolumntype{L}[1]{>{\raggedright\arraybackslash}p{#1}}
\setlength{\tabcolsep}{5pt}

\begin{document}

\vspace*{5mm}
\hrule
\vspace{3mm}
\begin{center}
    {\large\textbf{INSTITUTO DE COMPUTAÇÃO - UNICAMP}} \\
    \vspace{3mm}
    {\large\textbf{MO810A/MC959A - Sistemas Multiagentes com Modelos de Linguagem}} \\
    \vspace{3mm}
    {\large\textbf{Prof. Dr. Julio Cesar dos Reis}}
\end{center}
\vspace{3mm}
\hrule
\vspace{2.5cm}
\begin{flushright}
    {\Large\textbf{Fase 2 - Metodologia}} \\
    \vspace{1cm}
    {\Large\textit{EduAgent-OS: Sistema Multiagente Local para Tutoria Ativa e Gestão Longitudinal de Estudos}} \\
    \vspace{5cm}
    Rafael Attilio Agricola - 249245 \\
    Jociclelio Castro Macedo Junior - 231722
\end{flushright}

\newpage
\pagestyle{fancy}
\fancyhf{}
\renewcommand{\headrulewidth}{0pt}
\fancyfoot[R]{\thepage}

\section{Visão Geral da Metodologia}
\label{sec:visao_geral}

O EduAgent-OS investiga a integração de tutoria textual, memória entre sessões e planejamento de estudos para estudantes de graduação. O problema é manter um acompanhamento coerente quando materiais, dificuldades, prazos e disponibilidade mudam: responder uma pergunta não basta para determinar o que revisar nem quando fazê-lo. O projeto mantém o \textbf{ODS 4 -- Educação de Qualidade} como orientação principal, por buscar apoio à aprendizagem autônoma e à organização do estudo. A execução local também se relaciona aos ODS 9 e 10, mas a acessibilidade dependerá dos requisitos de hardware medidos. Não se pressupõe benefício educacional ou redução efetiva de custos antes da avaliação.

A metodologia compreende especificação de contratos, implementação incremental, construção de dados com referências independentes e experimentos controlados. O núcleo v0.1 processará PDFs e texto fornecidos pelo estudante, produzirá explicações fundamentadas, exercícios, flashcards e autoexplicações, registrará evidências formativas por tópico e proporá agendas semanais exportáveis em ICS. Curadoria web e sincronização remota serão capacidades opcionais, admitidas separadamente. O recorte é textual e inicialmente local, sem notas institucionais, treinamento fundacional, voz, geração de vídeos ou integração com sistemas acadêmicos fechados.

Há um ajuste operacional em relação à proposta: o estado de aprendizagem inicial será observado e auditável, com estimativa numérica nula. Um rótulo provisório de evidência favorável substituirá a inferência imediata de domínio. Da mesma forma, A2A será uma extensão possível para serviços distribuídos, enquanto a comunicação interna ocorrerá por artefatos tipados no grafo. Essas decisões tornam o primeiro protótipo verificável sem depender de parâmetros cognitivos não calibrados ou de distribuição de serviços.

\subsection{Hipóteses e ciclo integrado}

Serão investigadas três questões: (Q1) se a especialização em cinco agentes melhora a continuidade e o sucesso das tarefas frente a um agente único com recursos equivalentes; (Q2) se recuperação seletiva, compactação e seleção de habilidades reduzem custo preservando qualidade; e (Q3) se contratos, verificadores e limites contêm falhas antes que produzam efeitos indevidos. São hipóteses do projeto, não consequências garantidas pela literatura.

A organização por papéis e artefatos adapta elementos do MetaGPT e do MultiTutor \cite{hong2024metagpt,sun2025multitutor}. A separação entre contexto ativo e memória persistente se apoia no MemGPT \cite{packer2023memgpt}; habilidades reutilizáveis são inspiradas no Voyager \cite{wang2023voyager}. O emprego de ferramentas para operações verificáveis dialoga com o ToRA \cite{gou2024tora}, mas o solver de agenda e seus contratos são decisões próprias de engenharia.

O ciclo começa com a captura de materiais e restrições. O Curador extrai conteúdo e fatos candidatos, cuja confirmação antecede sua utilização como restrições duras. Para uma pergunta ou prática, recuperam-se fontes e memória do mesmo snapshot; o Tutor produz um turno ou item e encerra a invocação ao esperar o estudante. A resposta seguinte é avaliada pelo Monitor com referência e rubrica explícitas. Um redutor determinístico persiste a evidência e atualiza o estado por tópico. Quando houver dificuldade, tarefa perdida ou mudança confirmada, o Planejador propõe demandas e aciona um solver. Um verificador separado confere o plano antes do preview e da aprovação. A próxima sessão utiliza o estado atualizado, com origem e temporalidade preservadas.

\begin{figure}[htbp]
\centering
\begin{tikzpicture}[
    node distance=6mm and 8mm,
    box/.style={draw,rounded corners,align=center,font=\small,minimum height=10mm,text width=40mm,inner sep=3pt},
    wide/.style={box,text width=133mm},
    arr/.style={-{Stealth[length=2mm]},thick},
    infra/.style={box,fill=gray!8}]
\node[wide] (ui) {Estudante: materiais, perguntas, respostas e disponibilidade};
\node[wide,below=of ui] (g) {LangGraph: intenção, snapshot, contratos, orçamento e checkpoint};
\node[box,below=of g,xshift=-47mm] (c) {Curador\\Extração e RAG};
\node[box,below=of g] (b) {Bibliotecário\\Memória temporal};
\node[box,below=of g,xshift=47mm] (p) {Planejador\\Demandas e prioridades};
\node[box,below=of c] (t) {Tutor\\Explicação e prática};
\node[box,below=of b] (m) {Monitor\\Rubrica e evidências};
\node[infra,below=of p] (s) {CP-SAT + verificador\\Preview e ICS};
\node[wide,below=of m] (db) {Serviços determinísticos: validação, redutor e commit idempotente\\Documentos versionados, eventos/projeções SQLite e outbox};
\node[wide,below=of db] (llm) {Bonsai 2 27B local compartilhado: servidor único, fila e adaptador\\Política de ferramentas, cancelamento, limites e traces};
\draw[arr] (ui)--(g);
\draw[arr] (g.south west)--(c.north);
\draw[arr] (g)--(b);
\draw[arr] (g.south east)--(p.north);
\draw[arr] (c)--(t);
\draw[arr] (b.south west)--(t.north east);
\draw[arr] (t)--node[above,font=\scriptsize]{resposta}(m);
\draw[arr] (m)--(db);
\draw[arr] (p)--(s);
\draw[arr] (m.north east)--(p.south west);
\draw[arr] (s.south)--(db.north east);
\draw[arr] (db)--(llm);
\end{tikzpicture}
\caption{Arquitetura lógica proposta. As setas resumem dependências de dados; a resposta do estudante inicia uma nova invocação. Os cinco papéis compartilham a inferência local, e efeitos são executados por serviços determinísticos.}
\label{fig:arquitetura}
\end{figure}

Especialização significa instruções, entradas, permissões e critérios de término distintos, não cinco cópias do modelo nem diálogo irrestrito entre agentes. O grafo controla dependências, e nenhum parecer textual substitui a validação formal de um plano, o fato confirmado ou o recibo de persistência. O pacote documental das etapas 01--08 complementa a rastreabilidade; as regras necessárias à implementação e à avaliação estão descritas neste relatório. Todos os parâmetros quantitativos apresentados são propostas iniciais, sujeitas a piloto e congelamento anterior ao teste final.

\section{Conjunto de Dados para a Avaliação}
\label{sec:dados}

\subsection{Fontes, captura e registro}

O corpus será obtido de duas disciplinas de graduação autorizadas: uma de conteúdo técnico-formal e outra com respostas conceituais abertas. Incluirá PDDs, cronogramas, notas e materiais didáticos em português, com inglês técnico, PDFs digitais, documentos escaneados, tabelas e revisões conflitantes. A escolha busca diversidade de transformação e julgamento, sem representar estatisticamente todas as disciplinas. A obtenção dependerá de licença ou autorização documentada; materiais restritos não serão redistribuídos sem permissão.

Cada captura registrará identificador da disciplina, origem, data de obtenção, licença/autorização, MIME verificado pelos bytes, tamanho, SHA-256 e vínculo de versão. O original será preservado separadamente dos derivados. Fontes web opcionais serão congeladas em snapshots para comparações reproduzíveis, com URL final, estado de acesso, data e hash; consultas reais ao conector serão ensaios de integração separados. Um vídeo sem transcrição acessível oferecerá apenas metadados, não conteúdo supostamente lido.

Casos pessoais serão preferencialmente sintéticos: disponibilidade, compromissos e respostas roteirizadas não exigem históricos identificáveis de estudantes reais. Dados sensíveis presentes em documentos serão removidos ou substituídos antes da anotação e distribuição. Os exemplos públicos, manifestos e instruções de acesso autorizado permitirão conferir o corpus sem expor materiais restritos. Os bancos e caminhos de runtime serão segregados por usuário desde o esquema.

\subsection{Pré-processamento, qualidade e ontologia}

A ingestão utilizará Docling para blocos, ordem de leitura, tabelas e proveniência \cite{docling}. OCR seletivo em português ou português/inglês será aplicado quando a qualidade da extração justificar, preservando o original e o diagnóstico. Após OCR, o documento completo será reextraído: um sidecar isolado não representa necessariamente todo o texto \cite{ocrpdf}. Equações, células ou páginas ilegíveis permanecerão pendentes, sem preenchimento automático por inferência. Os limites iniciais são 50 MiB e 500 páginas por documento, 600 s por job e subprocesso limitado a 4 GiB; documentos maiores serão recusados ou divididos pelo usuário.

Duplicatas exatas poderão compartilhar bytes, mas manterão seus vínculos de disciplina. Traduções, scans, revisões e paráfrases serão registradas como parentes, não confundidas com fontes independentes. A normalização de avaliação usará NFC e espaços, preservando sinais, unidades e símbolos. A indexação lexical poderá usar NFKC e minúsculas, com tokenizer próprio versionado; essa transformação não alterará o original nem a transcrição de referência. Datas e números críticos serão comparados com o trecho de origem, não apenas com uma frase gerada.

A ontologia relacional terá entidades de usuário, disciplina, tópico, documento, versão, avaliação, recurso, item, tentativa, demanda e bloco de agenda. Tópicos manterão aliases e relações de pré-requisito; ciclos serão rejeitados. Um fato acadêmico conservará expressão original, valor normalizado, precisão, zona quando aplicável, fontes, autoridade e estado candidato, ambíguo, confirmado ou rejeitado. Dia sem horário não será convertido silenciosamente em instante. Datas relativas dependerão da data da fala e da zona confirmada. O PDD confirmado terá prioridade sobre sugestões, mas uma revisão documental gerará reconciliação, não substituição automática.

Os blocos serão transformados em chunks por seção, com alvo de 384 tokens do modelo de embeddings e sobreposição de até 64 quando necessária à continuidade. Tabelas repetirão cabeçalhos e conservarão referências de células/blocos. Cada referência de fonte identificará documento, versão, página física, bloco, offsets ou caixa e hash do texto. Uma nova versão só será ativada após a construção completa do índice; originais, blocos e índices manterão manifestos compatíveis.

\subsection{Composição, partições e prevenção de vazamento}

As metas iniciais são 24 famílias documentais, divididas em 8 para desenvolvimento, 4 para validação e 12 para teste; 120 consultas e 120 tentativas por partição; 30 trajetórias de teste com oito sessões; 100 agendas pequenas e 100 realistas; e pelo menos dez casos base para cada falha F01--F24. Agendas terão partições de 40/20/40\% por família de construção e parâmetros. Esses números são metas de coleta, não tamanhos já obtidos nem demonstração de potência suficiente.

Por disciplina, serão propostas 12 famílias (4/2/6 nas partições), com pelo menos dois scans e duas tabelas no corpus total. Cada split terá 80 consultas com suporte e 40 sem suporte, além de 60 tentativas por disciplina. Antes das 30 trajetórias de teste, serão construídas seis de desenvolvimento e quatro de validação, também com oito sessões e sem parentesco entre partições. Os autores selecionarão fontes e autorizações; um anotador preparará transcrição/ontologia, outro realizará exame independente dos originais, e um especialista da disciplina adjudicará discordâncias.

Consultas abrangerão suporte completo, parcial e ausente, literais técnicos, códigos, múltiplos aspectos e temporalidade. Tentativas incluirão respostas corretas, parciais, incorretas e indecidíveis, com diferentes níveis de ajuda e revelação de solução. Trajetórias conterão mudança de disponibilidade, erros recorrentes, correção temporal, revisão tardia e replanejamento. Scripts serão fixos, com alternativas aceitáveis e precondições explícitas; uma pergunta inesperada do sistema produzirá \code{SCRIPT_MISMATCH}, sem resposta favorável inventada pelo avaliador.

Toda transformação herdará a partição de sua família: tradução, scan, revisão, paráfrase, item derivado e variante adversarial não serão separados entre desenvolvimento e teste. Um auditor verificará hashes, similaridade, referências e o grafo documento--item--trajetória; candidatos a parentesco serão revisados por humanos, sem assumir limiar universal de similaridade. Gabaritos e rubricas privados do harness não entrarão no índice acessível ao sistema.

Os clusters de análise serão componentes de parentesco. Trajetórias sem fontes compartilhadas entre clusters serão preferidas. Se a conectividade reduzir muito o número de componentes, a análise usará esses componentes e explicitará a limitação, em vez de tratar 30 trajetórias ou centenas de turnos como estudantes independentes. Disciplina será um estrato descritivo, não uma unidade com apenas duas observações suficiente para generalização.

\subsection{Ouro, anotação e disponibilidade aos agentes}

O ouro será elaborado por dois especialistas independentes, com piloto em desenvolvimento e adjudicação de discordâncias por terceiro avaliador ou adjudicador. Serão preservados os julgamentos anteriores à adjudicação para medir acordo. O registro conterá autoria da anotação, versão, referências originais, fatos críticos, entidades/relações, relevâncias de recuperação, claims necessárias, respostas aceitas, critérios, códigos de erro, ledger temporal, restrições originais e efeitos esperados. Campos não aplicáveis ficarão ausentes, evitando interpretar listas vazias como negativos anotados.

Para extração, haverá transcrição humana e ouro de página, ordem, sinais, números e unidades. Para recuperação, qrels serão definidos sobre unidades originais e mapeados aos chunks, evitando que a configuração de segmentação determine o próprio ouro. Para itens, o especialista resolverá o enunciado antes de conferir o gabarito proposto; equivalências e tolerâncias serão explícitas. Cada critério da rubrica terá ID, \code{topic_ids} não vazio, níveis 0--2, marcação crítica, equivalências aceitas e referências de evidência. Critérios multicomponentes serão decompostos quando necessários à atribuição por tópico; sem atribuição separável, não haverá atualização desses tópicos. Para memória, o ledger distinguirá tempo de ocorrência, registro e validade. Agendas terão demandas e compromissos em instantes/minutos originais, com ótimo vetorial conhecido nas instâncias pequenas.

Em runtime, o agente receberá somente os dados autorizados da tarefa e do snapshot: fontes relevantes, fatos vigentes, item e rubrica quando cabíveis. O gabarito ficará separado da pergunta apresentada ao estudante. Os contratos \code{CaseRecord}, \code{GoldRecord} e \code{RunRecord} ligarão caso, família, cluster, partição, roteiro, snapshot, fontes, ouro e rubrica por IDs e hashes. Dados de teste permanecerão reservados ao runner após o congelamento. Vieses de cobertura linguística, qualidade de OCR, extensão de resposta e preferência de anotadores serão registrados por estrato, para não serem ocultados por uma média agregada.

\section{Arquitetura Multiagente}
\label{sec:arquitetura}

\subsection{Modelo compartilhado e admissão condicional}

Os cinco agentes utilizarão o mesmo candidato, \code{prism-ml/Ternary-Bonsai-2-27B-gguf}, revisão \nolinkurl{b072e1d3b35a0a630cece372c2127528e0994386}, com arquivo inicial \code{Ternary-Bonsai-2-27B-PTQ1_0.gguf} \cite{bonsai2}. O SHA-256 publicado do arquivo é \nolinkurl{53107f530aa52eb00912263ab1ee29bd199261c87cd7b4ad4ca1318c1fe33ee3}; a implementação deverá recalculá-lo localmente e conferir licença e NOTICE. A identificação evita que o nome Bonsai 2 seja tratado como configuração suficiente.

O runtime candidato será o fork \code{PrismML-Eng/llama.cpp}, release \code{prism-b10743-adfffbe}, fixada por hash após verificação \cite{bonsaidemo}. Não se presume que o llama.cpp padrão aceite o artefato ou possua as mesmas capacidades. Haverá um servidor loopback, um slot de geração e fila de quatro pedidos, sem visão. O manifesto registrará CPU, RAM, GPU/VRAM quando houver, sistema operacional, binários, quantização, tokenizer, template, offload e versões reais de dependências.

Antes de admitir o modelo, serão realizados carregamento offline, smoke do tokenizer e do template, contagem da janela, geração em português, parse de artefatos e rejeição de ferramentas proibidas. JSON Schema será solicitado ao servidor apenas se suportado pelo fork; validação estrutural e de política continuará obrigatória. O adaptador de geração devolverá texto final, uso, motivo de término e erro tipado. Falha de carregamento ou dos gates produzirá \code{MODEL_NOT_ADMITTED}; trocar de modelo exigirá revisão registrada, sem substituição silenciosa.

O perfil inicial sem pensamento explícito usará temperatura 0,7, top-p 0,8, top-k 20 e min-p 0. Em desenvolvimento será comparado ao perfil thinking com 1,0/0,95/20/0,05, conforme documentação do candidato. Template, penalidades e perfil selecionado serão congelados. Pensamento interno não será evidência pedagógica nem oráculo de avaliação; raciocínio e resposta, quando ambos gerados, compartilharão a reserva de saída. A escolha do modelo é motivada pelo objetivo local e pela possibilidade de identificar artefatos, não por eficácia ou latência já verificadas.

\subsection{Papéis, entradas, saídas e término}

\paragraph{Curador de Conhecimento.}
Receberá arquivos autorizados, metadados da disciplina, consultas e snapshot. Coordenará extração, diagnóstico, identificação de tópicos e fatos candidatos, recuperação documental e curadoria opcional. Produzirá manifesto/blocos, propostas de fatos e \code{EvidencePack} com chunks, ranking, fontes e cobertura suportada, parcial ou ausente. Pode propor um alias ou conflito, mas não confirmar prazo ambíguo nem transformar conteúdo documental em instrução. A tarefa termina em versão indexada, pacote válido, pendência explícita ou falha controlada; ingestão cancelada não ativa versão incompleta.

\paragraph{Tutor Adaptativo.}
Receberá objetivo, skill selecionada, fontes, memória pertinente e item pendente, sem o ledger privado do avaliador. Produzirá \code{TutorTurn} e, quando solicitado, \code{LearningItem} com pergunta, fontes, referência de resposta e rubrica separadas. Decidirá a forma de explicação e o próximo apoio pedagógico, sem alterar agenda ou declarar domínio. Seus estados serão explicar, tentativa, pista, autoexplicação, transferência e encerramento. Uma solicitação direta poderá receber explicação direta. Ao pedir resposta, termina em \code{waiting_student}; ao concluir, publicar solução solicitada, receber recusa ou cancelamento, encerra sem novas perguntas compulsórias.

\paragraph{Monitor de Aprendizagem.}
Receberá enunciado, resposta original, gabarito/rubrica, fontes e registro de ajuda/revelação, nunca apenas a autovalidação do Tutor. Produzirá \code{Assessment} por critério/tópico, tipo de evidência, código de misconception, corretude correta/parcial/incorreta/indecidível, confiabilidade aceita/abstenção e recomendação de revisão. Verificação exata, tolerância numérica e unidades prevalecerão quando aplicáveis. Texto aberto usará níveis 0, 1 e 2 com razões curtas. Evidência \code{present} exigirá offsets válidos na resposta; omissão será \code{evidence_kind=missing}, com \code{span=null} e descrição do critério ausente, podendo receber nível 0 sem abstenção automática. Span alegado mas inexistente será inválido; ambiguidade essencial ou divergência formal não resolvida exigirá abstenção. A tarefa termina em avaliação válida ou abstenção explícita; feedback útil poderá ser publicado sem atualização inferencial.

\paragraph{Planejador de Rotina e Agenda.}
Receberá fatos confirmados, disponibilidade, compromissos, demandas existentes, estado por tópico e pedidos de revisão. Produzirá \code{StudyDemand} e \code{ScheduleProposal} com status, blocos, déficit, diff, objetivo/bound e explicação. O LLM interpreta necessidades e prioridades; o solver decide a viabilidade da formalização e o verificador decide a admissibilidade da publicação. Duração recomendada é estimada até confirmação quando se tornar obrigatória. A tarefa termina em esclarecimento, preview viável, conflito provado para a grade ou resultado inconclusivo; aprovação é uma entrada posterior vinculada à revisão.

\paragraph{Bibliotecário de Memória.}
Receberá usuário, disciplina, tópico, intervalo temporal, snapshot e propostas de registros. Produzirá \code{MemoryPack}, sínteses com claims ligadas a eventos e propostas de correção, sem escrita irrestrita no banco. Distinguirá observado, declarado e inferido, atual e histórico, evidência independente e assistida. Não reativará fato substituído por similaridade textual. A tarefa termina em pacote resolvível, síntese verificada, contradição pendente ou ausência de evidência. A persistência final será executada pelo serviço determinístico.

\subsection{RAG, planejamento limitado e habilidades}

A recuperação será híbrida: BM25 local e BGE-M3 denso, com revisão/hash de embeddings fixados \cite{m3}. A busca lexical inicial usará \code{BM25Okapi}, $k_1=1{,}5$ e $b=0{,}75$; vetores normalizados e produto interno exato em arrays NumPy evitarão um serviço adicional no primeiro corpus. Usuário, disciplina e versão serão filtrados antes do ranking. Cada modalidade recuperará até 40 candidatos, combinados por
\begin{equation}
 \operatorname{RRF}(d)=\sum_{m\in\{lex,denso\}}\frac{1}{60+\operatorname{rank}_m(d)},
\end{equation}
com posições a partir de 1 e desempate pelo ID. A referência de RRF foi conferida por metadados; a operacionalização é também apoiada na documentação de busca híbrida \cite{rrf09,qdranthyb}. O reranker \code{BAAI/bge-reranker-v2-m3} será opcional, com pares de até 512 tokens e truncamentos registrados \cite{bgerank}. Sua ativação dependerá de ganho em validação e co-residência viável no hardware.

O contexto documental terá até seis chunks e 4.096 tokens do Bonsai. Toda citação resolverá a versão/página/bloco no snapshot. Resolução estrutural não prova apoio semântico; os dois critérios serão avaliados separadamente. Sem fonte suficiente, o sistema declarará a ausência; uma explicação geral poderá ser rotulada como tal, sem citação fictícia. O limiar de abstenção será escolhido em validação, sem converter escore de ranking em probabilidade calibrada.

O planejamento dos agentes será curto: selecionar rota e skill, gerar proposta estruturada, consultar ferramenta permitida, validar e publicar ou reparar dentro do orçamento. Não haverá autoedição de política ou conversação multiagente aberta. As skills versionadas terão ID, tarefas, schemas de entrada/saída, pré-condições, passos, ferramentas permitidas, pós-condições, término, máximo de chamadas e casos de referência. O catálogo obrigatório conterá explicar, pistas graduais, quiz, flashcard e autoexplicação; mudanças ocorrerão em revisão offline e terão hash registrado.

A prática de recuperação e a autoexplicação são motivadas por trabalhos pedagógicos \cite{roediger2006testing,chi1994selfexplanation}, sem transportar seus efeitos automaticamente para este protótipo. As pistas terão três níveis: conceito, próximo passo e exemplo análogo. Após três, o Tutor oferecerá solução comentada ou encerramento; solução solicitada e recusa interromperão a sequência. Flashcards terão conceito atômico, frente/verso separados e fonte. Intervalos de revisão de 1/3/7 dias serão política inicial configurável. Transferência utilizará famílias identificadas, impedindo que cópia de solução conte como evidência independente.

\subsection{Memória temporal, compactação e estado observado}

SQLite armazenará eventos originais e projeções versionadas. Eventos terão \code{occurred_at} e \code{recorded_at}; fatos terão intervalo de validade, registro e relação \code{supersedes}. Preferência declarada não equivale a evidência de aprendizagem. Consulta atual utilizará projeção vigente e correções; consulta histórica especificará se o corte é por validade ou pelo que já estava registrado naquele instante. Contradição não conciliada manterá ambos os registros e a pendência.

A recuperação de memória filtrará usuário, disciplina, tópico e tempo antes da busca lexical/densa, selecionando até oito originais em 2.048 tokens, além de fatos atuais necessários. Correções relevantes terão prioridade sobre vizinhos semelhantes. Essa política será testada com consultas multissessão, atualizações e abstenção, adaptando categorias do LongMemEval \cite{wu2025longmemeval}.

O contexto será contado com o tokenizer real após template e serialização. A janela operacional proposta é
\begin{equation}
 T_{\mathrm{entrada}}+16\,384+2\,048+1\,024\leq32\,768,
 \qquad T_{\mathrm{entrada}}\leq13\,312.
\end{equation}
As reservas correspondem a saída, ferramentas e margem; limites locais de fontes e memória não autorizam ultrapassar o total. A montagem seguirá instrução/schema, skill, tarefa/item/rubrica, âncoras confirmadas, fontes, memória e turnos recentes. Antes de resumir, serão removidas fontes periféricas ou segmentada a tarefa. Rubrica indispensável não poderá ser descartada para continuar avaliando.

Sínteses serão derivadas dos originais, com claim--IDs de eventos e cobertura temporal. Datas, números, polaridade, ajuda e estado serão conferidos explicitamente. Síntese inválida não substituirá o histórico: conservar-se-ão originais selecionados ou será solicitada divisão. Não haverá cadeia indefinida de resumo de resumo, e compactar não apagará eventos. A ablação sem compactação usará a mesma seleção de originais e o mesmo limite; falta de espaço será registrada como impedimento, não compensada com janela maior.

O redutor utilizará, para cada tentativa, o julgamento vigente: a maior \code{assessment_revision} explicitamente aceita, com \code{supersedes} e evento de correção. A unicidade lógica será \code{(attempt_id,assessment_revision)}. Retry não duplicará efeito; revisão de julgamento substituirá a contribuição anterior por replay dos eventos, sem somar uma segunda atualização cumulativa. Cada tentativa terá \code{occurred_at} e \code{attempt_ordinal} atribuído no commit; a ordem será por ocorrência, com desempate pelo ordinal. Julgamento tardio ou corrigido recalculará a sequência sem mudar a ordem da tentativa. Independência exigirá pista de nível 0, ausência de reveal e de acesso prévio à solução da família na sessão corrente; exposição em outra sessão será marcada como revisão, não como ajuda atual automática.

A corretude de um tópico considerará somente seus critérios: todos em nível 2 serão corretos, todos em 0 incorretos e mistura parcial; critério essencial indecidível tornará o tópico indecidível, sem atualização de competência. Critério crítico com nível 0 definirá erro crítico. A estimativa será \code{null}. Sem evidência, o estado será \code{insufficient}; nos demais casos, \code{mixed}, salvo a regra de \code{favorable}: as três tentativas independentes, aceitas e decidíveis mais recentes do tópico devem ser todas corretas, de famílias distintas e posteriores ao último erro crítico aceito, considerando os julgamentos vigentes. ``Mais recentes'' designa ordem de tentativa, não janela de dias; resposta parcial não conta como correta.

Um novo erro crítico aceito, mesmo assistido, retorna a \code{mixed} e solicita revisão. Uma nova tentativa independente parcial ou incorreta também interrompe a sequência favorável. Acertos assistidos não removem erro crítico nem compõem a sequência; leitura e autorrelato não alteram competência. Evidência antiga será indicada pela idade, sem inventar erro por ausência de prática. O rótulo favorável é provisório e não significa domínio. BKT será apenas extensão futura em shadow mode, ajustada em logs separados e admitida após calibração prospectiva \cite{badrinath2021pybkt}.

\subsection{Política de evidências para demandas}

O serviço determinístico aplicará \code{assessment_review_policy_v1}, registrada em demandas e \code{review_due}. Tentativa aceita parcial/incorreta ou erro crítico gerará revisão sugerida no próximo dia local com disponibilidade; na falta dele, buscará a primeira data disponível no horizonte. Acerto independente gerará candidatos de recuperação em 1/3/7 dias a partir da ocorrência local. A relação \code{assessment_topics} ligará tópicos a avaliações acadêmicas confirmadas; o prazo pedagógico será o menor entre o due sugerido e o início confirmado da avaliação. Avaliação informada apenas por dia motivará sugestão anterior à data, sem criar datetime duro. Revisão vencida conservará o due original e buscará bloco futuro com atraso indicado; além do horizonte, ficará pendente para o seguinte.

Prioridade opcional de classe 1 corresponderá a avaliação em até 48 h ou erro crítico; classe 2, avaliação em até sete dias, revisão vencida ou dificuldade recorrente (dois erros aceitos de famílias distintas nas últimas cinco tentativas do tópico); classe 3, demais revisões. Sem avaliação conhecida, não se inventará urgência. Duração informada/confirmada prevalecerá; na ausência, o catálogo proporá \code{practice=30min}, \code{review=15min} e \code{flashcard=15min}, com \code{estimated} e \code{mandatory=false}. Edição e confirmação do estudante serão necessárias para promover a demanda a obrigatória; percentual de domínio não determinará duração.

Haverá no máximo uma demanda pendente por usuário, tópico, tipo de atividade, due local e versão da política. Consolidação reunirá todas as \code{evidence_ids} e a prioridade mais alta, sem duplicar demandas em retry. \code{review_due} será a data do primeiro candidato pendente, não probabilidade de esquecimento. O Planejador poderá explicar ou propor alteração, mas o serviço aplicará a política e exigirá confirmação para mudar restrição dura. A avaliação distinguirá aderência à política e pertinência pedagógica da recomendação.

\subsection{Formalização do planejamento temporal}

O horizonte irá da meia-noite local da data inicial confirmada à meia-noite após sete datas locais, convertendo cada fronteira para UTC. A grade de 15 minutos nascerá às 00:00 de cada dia local; candidatos não atravessarão meia-noite no MVP. UTC servirá para comparações e uma zona IANA confirmada para janelas locais. Disponibilidade será arredondada para dentro e bloqueios para fora; uma demanda de $d_i$ minutos ocupará $\lceil d_i/15\rceil$ slots, sem reduzir duração. Dia sem hora e horário local ambíguo/inexistente exigirão esclarecimento. Capacidade diária contará minutos efetivamente reservados, inclusive arredondamento e blocos fixos de estudo; compromissos pessoais bloquearão tempo sem consumir a capacidade de estudo. Fixos serão inseridos uma única vez na ocupação, sem candidatos móveis.

Para cada demanda $i$, $B_i$ será o conjunto de candidatos contíguos compatíveis com início mínimo, prazo e janelas. Com $x_b\in\{0,1\}$ e presença $p_i=\sum_{b\in B_i}x_b$, o CP-SAT imporá
\begin{align}
 p_i&=1 &&\text{se $i$ obrigatória},\\
 p_i&\leq1 &&\text{se $i$ opcional},\\
 \sum_{b:t\in b}x_b&\leq1 &&\text{para cada slot $t$},\\
 p_j&\leq p_i &&\text{se $i$ é predecessor pendente confirmado de $j$}.
\end{align}
Se ambas as demandas pendentes de uma precedência estiverem presentes, $end_i\leq start_j$. Dependências temporais referenciarão demandas/atividades, não IDs de tópicos: pré-requisito pedagógico só se tornará dependência de atividade quando explicitado e confirmado. Predecessor concluído satisfará presença com término real; predecessor iniciado ocupará intervalo fixo e o dependente aguardará término confirmado. Predecessor pendente fora do horizonte bloqueará o dependente com razão explícita. Serão impostos limites diários confirmados e imutabilidade dos fixos; ciclo será erro de entrada. Tarefas serão indivisíveis por padrão. Decomposição exigirá \code{splittable}, número/duração de subtarefas e soma pelo menos igual à duração original, confirmados pelo usuário. A documentação de CP-SAT orienta status e implementação, não valida a interpretação das demandas \cite{google2024cpsat}.

O objetivo será lexicográfico: maximizar atendimento opcional de classe 1, depois 2 e 3; minimizar blocos futuros previamente aprovados removidos ou com instantes alterados; minimizar deslocamento em slots, somando diferenças absolutas de início em minutos divididas por 15 sobre IDs mantidos; e maximizar overlap em minutos com janelas preferidas declaradas. Adições serão contadas separadamente; ausência de preferência terá valor zero. Os dez segundos totais serão divididos por tempo restante/etapas restantes, reaproveitando sobra. Cada etapa conservará incumbent validado e vetor de provas; objetivo só será fixado após \code{OPTIMAL}. \code{FEASIBLE} interromperá a lexicografia. \code{UNKNOWN} posterior a uma etapa ótima devolverá o incumbent anterior como \code{FEASIBLE}, com provas parciais, sujeito ao verificador independente; sem qualquer solução, será inconclusivo. \code{INFEASIBLE} será apresentado como inviável na grade conservadora de 15 minutos, com viabilidade contínua não determinada; \code{MODEL_INVALID} será defeito.

Um verificador implementado separadamente reaplicará as restrições em minutos e instantes originais, conferindo duração, conflitos, precedência, demandas, imutabilidade e snapshot. Relaxação diagnóstica nunca será publicada como plano aprovado. O preview mostrará adições, movimentos, remoções, demandas não atendidas e origem das prioridades. A aprovação vinculará ID, revisão, digest e dependências à revisão global do domínio; sua validade para outbox e envio seguirá a política da Seção~\ref{sec:orquestracao}.

\subsection{Registro de atividades e aprovação}

\code{ActivityUpdate} terá bloco, status, ocorrência, revisão esperada e ID de operação. As transições serão planejado--iniciado--concluído ou planejado--perdido/cancelado; conclusão direta de planejado exigirá início/fim reais informados. Sair de concluído exigirá correção explícita. \code{missed} será declaração do estudante ou sugestão de atraso confirmada, nunca inferência automática pelo relógio. Iniciado bloqueará o intervalo aprovado; extensão além do fim previsto exigirá atualização de término antes do solver e manterá conflito até esclarecimento. \code{approve_plan} validará o preview atual. Aprovação e atualização de atividade retornarão recibos e poderão disparar replanejamento uma vez por evento.

\section{Orquestração e Uso de Ferramentas}
\label{sec:orquestracao}

\subsection{Grafo, contratos e terminais}

O LangGraph coordenará rotas fixas, checkpoints persistentes e pausas \cite{lgstate,lginterrupt}. O estado terá usuário, sessão, turno, tarefa, trace, intenção, snapshot, referências de artefatos, pergunta/item pendente, orçamentos, deadline, cancelamento, status e erros. Checkpoints guardarão referências, não PDFs integrais. Leituras independentes do Curador e Bibliotecário poderão compartilhar snapshot; gerações seguirão a fila serial. A base de domínio será autoritativa, e checkpoint não substituirá commit.

\code{users.domain_revision} será o contador monotônico global por usuário: operação nova que mudar domínio o incrementará na mesma transação; retry, recibo, lease e trace não o incrementarão. \code{snapshot_version} e \code{expected_version} referenciarão esse contador. O snapshot será materializado em transação curta de leitura, copiando fatos, agenda, projeções, versões ativas e hashes necessários para artefato imutável com ownership. Agentes não manterão conexão aberta nem hidratarão silenciosamente dados atuais; índices consultarão exatamente as versões listadas. Nova escrita após a leitura exigirá recarga antes do próximo commit, mesmo quando a alteração não estiver relacionada à tarefa, conforme controle conservador do MVP.

As intenções serão ingestão, consulta, prática, submissão de resposta, planejamento, replanejamento, aprovação de plano, atualização de atividade, consulta/correção/exclusão de memória, exportação, curadoria e cancelamento. Comando explícito prevalecerá sobre classificação; classificador LLM só retornará enum validado. Entrada ambígua pedirá esclarecimento antes de efeito. Consulta passará por recuperação, contexto, Tutor, validação e publicação/espera; submissão resolverá item pendente, Monitor, validação, commit e feedback; agenda passará por demandas, solver, verificador e preview, com exportação apenas da revisão aprovada.

Contratos Pydantic estritos terão campos extras proibidos, enums, IDs UUID, referências ID+versão e tamanhos limitados \cite{pydstrict}. Datetimes serão ISO-8601 com zona/UTC; datas parciais terão tipo próprio. O envelope conterá \code{schema_version}, \code{task_id}, \code{trace_id}, \code{parent_task_id}, \code{sender_role}, \code{artifact_type}, \code{snapshot_version}, \code{source_refs}, \code{payload_ref}, \code{status} e erro tipado. Referências produzidas pelo LLM não serão caminhos de arquivo. O receptor conferirá papel, tipo, versão, ownership e resolução antes de usar o payload.

Terminais observáveis serão concluído, espera legítima, abstenção, conflito, erro, cancelado e efeito desconhecido. Espera pelo estudante encerra a invocação e não consome tempo ativo; uma nova entrada retoma o item. O reinício de um nó após interrupt exige preparação pura e efeitos em nós idempotentes separados. Divergência com fonte confirmada ou solver não será decidida por votação de agentes: haverá esclarecimento ou falha controlada.

\subsection{Matriz de ferramentas e validação}

Ferramentas serão expostas por adaptadores com allowlist; MCP poderá padronizar fronteiras externas \cite{mcparch}. Nenhum agente terá shell genérico ou escrita direta no banco. A autorização será decidida por código, mesmo que uma chamada proibida esteja em JSON válido. Documento, site e retorno de ferramenta serão dados rotulados, sem poder alterar papel ou política.

\begin{table}[htbp]
\centering\small
\caption{Ferramentas propostas; ``serviço'' identifica executor determinístico.}
\begin{tabular}{@{}L{3.3cm}L{2.5cm}L{4.0cm}L{5.0cm}@{}}
\toprule
Ferramenta & Proponente & Condição & Validação/resultado \\
\midrule
\code{document.extract} & Curador & Arquivo autorizado e job & MIME, limites, diagnóstico, página/bloco \\
\code{document.search} & Curador e Tutor & Disciplina e snapshot & Ownership, filtros, fontes e cobertura \\
\code{memory.query} & Bibliotecário & Usuário, tópico e tempo & Versões, correções e origem \\
\code{assessment.check} & Monitor & Item e rubrica válidos & Unidades, tolerâncias, spans ou ausência tipada \\
\code{schedule.solve} & Planejador & Restrições duras confirmadas & Status, bound e verificador independente \\
\code{calendar.render} & Planejador & Plano válido/aprovado para exportar & Parser distinto e equivalência \\
\code{web.search/fetch} & Curador & Rede/conector habilitados & Destino, quotas, conteúdo e acesso \\
\code{domain.commit} & Serviço & Proposta válida e revisão esperada & Transação, deduplicação e recibo \\
\code{calendar.apply} & Serviço & Confirmação atual e outbox & ID remoto, precondição e reconciliação \\
\bottomrule
\end{tabular}
\label{tab:tools}
\end{table}

A curadoria conectada terá até cinco candidatos e três fetches, 10 s e 5 MiB por recurso, até dois redirecionamentos revalidados. Guardará estados metadados, texto lido, transcrição lida e indisponível. Sem conector configurado, retornará \code{CONNECTOR_DISABLED}; cache offline será datado. Paths canônicos, ownership e bloqueio de traversal/symlink escapes protegerão arquivos. Fetch bloqueará destinos privados/loopback e revalidará DNS e redirects. Essas validações adaptam práticas de fronteira do MCP \cite{mcpsec}; sua efetividade será ensaiada por efeitos e canários, não presumida.

\subsection{Orçamentos, erros e intervenção do estudante}

Os limites propostos por turno são seis chamadas LLM, oito chamadas de ferramentas, 16 passos de grafo, dois reparos no total e um por artefato, 120 s ativos e 90 s por geração. Reparos consomem chamada, tempo e contador. Ingestão será job separado com progresso e cancelamento. Antes de cada nó, ferramenta ou geração, serão conferidos budget, deadline e cancelamento. Fila cheia retornará \code{RESOURCE_LIMIT}; esgotamento terminará com erro e artefatos já válidos, sem promover saída truncada.

Erros tipados incluirão schema inválido, fonte não resolvida, entrada ambígua, conflito de versão, ferramenta negada, modelo não admitido, ausência de evidência, limite de recurso, timeout, efeito desconhecido, cancelamento e falha interna. Cada erro terá mensagem, possibilidade de retry e referência de diagnóstico local. Cancelar impedirá novos despachos e reconciliará efeitos já iniciados; não apagará a necessidade de conhecer o estado remoto.

O estudante confirmará fatos ambíguos, prazos sem hora, restrições obrigatórias estimadas, revisão de plano e efeitos externos. Correção/exclusão mostrará escopo antes do commit. A confirmação não será genérica: ficará ligada ao artefato, revisão e digest. Sugestão de revisão poderá ser opcional enquanto sua duração não for confirmada, evitando transformar incerteza pedagógica em obrigação temporal.

\subsection{Transações, idempotência e outbox}

O armazenamento usará SQLite com chaves estrangeiras e constraints habilitados, WAL/FULL, escritor serializado e transações curtas \cite{sqlitewal}. A versão de runtime deverá incorporar as correções WAL documentadas, inicialmente 3.51.3 ou backport corrigido. Tabelas pessoais terão ownership direto ou herdado por FK composta \code{(user_id,id)} para o pai; eventos e itens pessoais incluirão usuário. Arestas de tópicos exigirão mesmo dono e disciplina. Originais/eventos serão separados de fatos, avaliações, estados, resumos, demandas, planos, operações, confirmações e outbox. Não haverá transação aberta durante inferência ou rede.

O commit receberá \code{operation_id}, digest, usuário e \code{expected_version}. Após conferir identidade/permissão, verificará \textbf{primeiro a identidade da operação}: mesmo ID, digest e usuário devolverão o recibo anterior, mesmo se a revisão esperada agora for obsoleta. Mesmo ID com digest ou usuário diferente será negado. Somente para operação nova será conferida a revisão global esperada; evento, operação, projeção, incremento de \code{domain_revision} e eventual outbox serão gravados atomicamente. Esse ordenamento evita rejeitar como conflito um retry de operação já efetivada.

Preview ficará ligado à revisão global e ao digest. A confirmação guardará \code{base_revision}, \code{approval_revision} pós-commit e conjunto/hash das dependências. Criar nova outbox exigirá domínio igual à revisão da aprovação; confirmação/outbox poderão compartilhar o commit. Antes do envio, o dispatcher distinguirá operação já confirmada de nova proposta e verificará que nenhuma restrição ou atividade dependente mudou. Mudança de dependência invalidará aprovação ainda não aplicada; efeito já aplicado exigirá novo diff e confirmação, sem undo silencioso. Metadados operacionais não invalidarão por si mesmos a revisão de domínio.

O checkpoint ficará em banco separado e não será atomicamente confirmado com o domínio. Após crash, replay consultará a operação autoritativa e reutilizará o recibo. Dispatcher com lease processará outbox apenas após confirmação atual. Antes de commit, crash não promoverá estado; após commit, recuperará pendência; após efeito e antes de ACK, consultará o destino pela identidade e reconciliará. Sem consulta possível, manterá \code{UNKNOWN_EFFECT}, sem repetir cegamente. A entrega poderá ser ao menos uma vez; não se promete exactly-once distribuído.

\code{artifact_dependencies} registrará usuário, referências de filho/pai e relações de derivação, correção ou substituição, indexadas nos dois sentidos. Correção/exclusão percorrerá dependentes, invalidará versões publicadas e reconstituirá projeções. Fonte compartilhada terá vínculo por usuário: excluir um vínculo não removerá bytes que outro dono ainda mantenha autorizados. Exclusão ativa será imediata na aplicação e propagará a payloads, índices, resumos, checkpoints e traces pessoais.

A retenção proposta de backups será de 30 dias; cópias com dados excluídos serão declaradas pendentes até expiração. Backup contendo apenas o usuário excluído será removido já na solicitação. Tombstones sem payload impedirão ressurreição e serão aplicados antes de liberar restore ao sistema, sem manter hash reversível do texto apagado. CLI mostrará escopo lógico e prazo, sem afirmar sanitização forense; job ou doctor verificará expiração com relógio e catálogo. Backup inicial fechará escritores/checkpointer, consolidará o banco e copiará conjunto consistente com originais e manifesto. Restore ocorrerá em diretório novo e passará por integridade e replay.

\subsection{ICS, sincronização e observabilidade}

A exportação seguirá RFC 5545, com biblioteca \code{icalendar}, VCALENDAR VERSION 2.0/PRODID e VEVENT com UID, DTSTAMP, DTSTART, DTEND e SUMMARY \cite{desruisseaux2009icalendar}. UTF-8, CRLF e folding serão validados por parser independente. Início/fim usarão UTC Z e fim exclusivo; all-day usará DATE e término no dia seguinte. A exportação simples será \textbf{sem METHOD}. \code{activity_occurrence_id} será persistido ao criar demanda: a mesma ocorrência conservará \code{block_id} e UID ao mudar horário; outra ocorrência receberá nova identidade. O bloco lógico armazenará \code{created_at}, \code{modified_at}, \code{event_sequence}, resumo, tipo de valor, demanda de origem e início/fim reais.

CREATED conservará a criação da ocorrência; DTSTAMP e LAST-MODIFIED representarão a modificação persistida da revisão do evento. SEQUENCE refletirá \code{event_sequence}, iniciada em zero e incrementada somente quando campos do evento mudarem, não por mudança em outra parte do plano. Retry e reexportação do mesmo evento não alterarão timestamps/SEQUENCE. Diff usará identidade de ocorrência; remoção na exportação ICS significará ausência no arquivo, não comando remoto de delete, que exigirá operação confirmada. UID não garante deduplicação em todo importador manual.

Sincronização será admitida apenas se o conector suportar leitura, identidade e precondições de versão. CalDAV utilizará href estável, \code{If-None-Match:*} para criação e \code{If-Match} com ETag para atualização \cite{daboo2007caldav}. Outros backends precisarão de verificação própria de IDs e condições de update/delete. Alteração externa exigirá conflito e novo preview. Conector sem garantias requeridas permanecerá não admitido, preservando preview e ICS locais.

Traces locais registrarão spans de fila, inferência, retrieval, contexto, solver, commit e ferramentas, com IDs, versões, seed, sampling, contagens, término e erros \cite{oteltrace}. Tokens incluirão todas as chamadas e reparos; raciocínio não será somado duas vezes quando já incluído na saída. Tempo ativo, fila e total serão separados. RSS/VRAM serão amostradas inicialmente a cada 100 ms, indicando limitação de picos. Payload pessoal completo não será logado por padrão. Métrica indisponível será nula com motivo, nunca zero. Energia e custo monetário dependerão de medidor ou amortização documentados, separando custo do sistema e custo da avaliação.

\subsection{Sequência de implementação}

Python 3.11+, lock e manifesto fixarão LangGraph, Pydantic v2, Docling/OCR, embeddings, OR-Tools e bibliotecas de calendário. A entrega terá CLI/API loopback, schemas exportados, migrações, skills e módulos de serviços; dados e checkpoints ficarão separados. O trabalho seguirá dependências verificáveis:
\begin{enumerate}
\item Admitir hardware/modelo, tokenizer, adaptador e instalação offline; iniciar registro de casos e auditoria de splits.
\item Implementar contratos, repositório, migrações, commit idempotente e recuperação antes de efeitos remotos.
\item Implementar ingestão, ontologia, proveniência e confirmação; comparar com ouro documental.
\item Implementar RAG híbrido e abstenção; selecionar reranker em validação.
\item Implementar skills, Tutor, Monitor e redutor; validar itens, rubricas e ajuda.
\item Implementar memória bitemporal, montagem/compactação e propagação de correção/exclusão.
\item Implementar demandas, CP-SAT, verificador separado, preview, confirmação e ICS.
\item Integrar grafo, pausas, cancelamento, checkpoints e traces; admitir conectores individualmente.
\item Concluir piloto e congelar dados, condições, métricas, budgets e hipóteses; executar o testbed e emitir resultados somente após as runs.
\end{enumerate}

\section{Procedimentos de Avaliação}
\label{sec:avaliacao}

\subsection{Testbed, famílias e oráculos}

A avaliação será em camadas: contratos/componentes determinísticos; qualidade dos 18 tipos de saída; fluxos W01--W12 completos; continuidade longitudinal; robustez F01--F24; e custo/reprodução. A cobertura integral por saída e falha consta do Apêndice. Instâncias pequenas terão referência formal independente; texto e pedagogia terão rubricas humanas. Métodos de ranking, citações e fidelidade adaptam BEIR, ALCE e RAGAs \cite{thakur2021beir,gao2023alce,es2024ragas}, sem tratar pontuação automática como prova de correção.

O \code{SUTAdapter} terá inicialização por manifesto/condição/sandbox/relógio/seed, carga de snapshot, invoke, resume, cancel, leitura de artefato, inspeção de domínio, exportação de trace, restart e close. O retorno conterá terminal, fontes, artefatos, operações, uso e erros. Inspeção autoritativa será recurso do harness, não ferramenta adicional do agente. MockLLM servirá aos testes determinísticos, sem ser braço de comparação de qualidade real.

O registro de execução terá par, condição, seed, braço limpo/perturbado, hashes, snapshot inicial, tempos, conclusão, terminal seguro, elegibilidade, trigger alcançado, falha injetada e erro do harness. Métricas manterão numerador, denominador, unidade, oráculo/versão, cluster e motivo de ausência. Erro do harness será separado de falha do sistema, sem substituir silenciosamente a run; repetição técnica seguirá regra pré-definida.

O registro ligará requisitos C01--C20 a casos e predicados executáveis. Validator rejeitará ID desconhecido e requisito do núcleo sem caso positivo, negativo e predicado. A matriz concreta será gerada após captura, com IDs reais de casos. Oráculos próprios reproduzirão revisão global/materialização, ownership direto ou composto, julgamento vigente e ordinal, omissão válida, política de demanda, atividades/dependências, grade e identidade de ocorrência. Capacidade desabilitada será N/A com motivo, sem crédito de cobertura por outra rota.

Os oráculos serão testados em desenvolvimento por controles válidos e mutantes conhecidos: sobreposição de um minuto, duração de 17 reduzida a 15, pré-requisito omitido, fixo movido, página existente mas errada, span inexistente, ajuda apagada, DTSTAMP inconsistente por revisão, canário ressurgido e contador sem retry. Um avaliador incapaz de detectar seu mutante será corrigido antes do congelamento. Serializer e parser distintos evitarão round-trip tautológico; incapacidade do parser de interpretar propriedade válida será marcada como interoperabilidade inconclusiva, não falha automática do serializer.

\subsection{Métricas de artefato e critérios propostos}

Extração medirá distância de edição de caracteres/palavras, $CER/WER=(S+D+I)/N$, por transcrição humana \cite{jiwerQuality}; ambas podem exceder 1. Serão reportados exact match de sinais, datas, números/unidades e ordem de blocos. Entidades/relações terão pareamento um a um, precisão $P=TP/(TP+FP)$, recall $R=TP/(TP+FN)$ e $F_1=2PR/(P+R)$; datas incluirão tipo, precisão, zona, autoridade e confirmação. Referência existente com página errada não passará no predicado de proveniência.

Para recuperação respondível, $Recall@40=|R_q\cap D_{40}|/|R_q|$ e $MRR@6$ será a média do inverso da posição do primeiro relevante até seis, ou zero quando ausente. Com ganho linear,
\begin{equation}
 DCG@6=\sum_{r=1}^{6}\frac{rel_r}{\log_2(r+1)},
 \qquad nDCG@6=\frac{DCG@6}{IDCG@6}.
\end{equation}
Query respondível sem ranking contará zero; query sem ouro relevante pertencerá ao estrato de abstenção, sem recall fictício. Completude de aspectos será anotada para perguntas multitrecho.

Precisão de citação (CP) contará citações em claims apoiadas pelo conjunto citado, em que cada citação apoia sozinha ou é necessária ao apoio conjunto; cobertura de citação (CR) contará claims documentais necessárias com citação que sustenta:
\begin{align}
 CP&=\frac{\#\text{citações adequadas}}{\#\text{citações avaliadas}},\\
 CR&=\frac{\#\text{claims necessárias sustentadas e citadas}}{\#\text{claims documentais necessárias avaliadas}}.
\end{align}
Fidelidade ao contexto, corretude factual e citação serão medidas distintas. Denominador zero será N/A com razão. Claims atômicas e fontes serão anotadas antes de utilizar um juiz auxiliar.

Explicações e turnos receberão rubrica 0--4 por corretude, pertinência, clareza, explicitação de limites e apoio contingente. Itens serão conferidos por especialista e, quando possível, verificação numérica $|a-a^*|\leq atol+rtol|a^*|$, unidades e equivalência simbólica sob hipóteses explícitas. O Monitor terá kappa ponderado linear por nível,
\begin{equation}
 \kappa_w=1-\frac{\sum_{a,b}d_{ab}O_{ab}}{\sum_{a,b}d_{ab}E_{ab}},
 \qquad d_{ab}=\frac{|a-b|}{2},
\end{equation}
além de macro-$F_1$ de códigos de erro, falso crédito e omissão de erro crítico. Kappa pode ser negativo ou indefinido. Falso crédito será reportado entre avaliações aceitas, juntamente com cobertura de aceitação e curva risco--cobertura, para não premiar abstenção universal. Avaliação por rubrica é motivada pelo Prometheus, sem supor que o Bonsai possua suas capacidades \cite{kim2024prometheus}.

Estado de aprendizagem será comparado por replay de prefixos com redutor de referência; memória terá acurácia atual, retrospectiva, multissessão, atualização e abstenção. Compactação medirá recall de tuplas críticas (sujeito, predicado, valor, polaridade, tempo, ajuda e fontes), contradições, tokens e qualidade downstream. Não serão calculados Brier ou log loss para \code{estimate=null}; se um módulo probabilístico futuro for admitido, $Brier=n^{-1}\sum_{i=1}^{n}(p_i-y_i)^2$ e log loss de próxima resposta exigirão calibração prospectiva \cite{guo2017calibration}. Log loss não tem teto e infinitude será registrada por flag, não JSON inválido.

Demandas serão comparadas às restrições originais confirmadas, incluindo omissões e duras extras. Uma implementação simples e independente da política v1 conferirá due, prioridade, duração estimada/confirmada, dedup e evidências, separadamente da pertinência humana. Planos terão violações por publicação, atendimento, déficit, status/provas parciais, capacidade em minutos reservados e churn por ocorrência. Nas pequenas agendas, enumeração própria verificará viabilidade e vetor objetivo; não se imporá identidade de plano em empates. ICS e sincronização terão equivalência por tuplas plano--arquivo--readback, UID e sequência do evento, incluindo plano alterado sem mudança do evento. Recibos terão ledger externo de efeito, atomicidade, invalidação de aprovação e exclusão propagada no escopo lógico declarado; backups pendentes serão conferidos até expiração e tombstones antes de restore.

Os gates iniciais propostos são: exact match crítico de extração $\geq0{,}98$; $F_1$ ontológico $\geq0{,}90$; recall de recuperação $\geq0{,}90$ e CP $\geq0{,}95$; mediana humana $\geq3$ por dimensão; validade de itens/cards $\geq0{,}95$; kappa do Monitor $\geq0{,}70$ e falso crédito $\leq0{,}05$ com cobertura; acurácia de memória $\geq0{,}90$; pertinência de curadoria $\geq0{,}90$. Proveniência, replay, fatos críticos de compactação, fidelidade de restrições duras, equivalência ICS, terminais e spans obrigatórios de trace exigirão 100\% no ensaio. Evidências present exigirão offsets válidos; missing exigirá span nulo e omissão descrita. Nenhum caso aceito poderá conter erro crítico; não se tolerarão plano inválido publicado, vazamento, efeito duplicado ou perda de dado comprometido. São critérios de engenharia para o corpus escolhido, não mínimos universais publicados nem garantia além da amostra. ICs e denominadores acompanharão todo gate; não será reduzido um limite depois de observar teste.

\subsection{Sucesso, disponibilidade, custo e delegação}

Uma tarefa terá sucesso somente se todos os predicados essenciais da rota e o terminal correto forem satisfeitos. Abstenção correta em caso sem suporte e conflito correto em agenda inviável poderão ser sucesso; erro controlado após falha induzida pode ser contenção segura sem concluir a tarefa original. Serão mantidas duas variáveis: \code{task_completion} e \code{safe_terminal}.
\begin{align}
 TSR&=\frac{\#\text{tarefas com sucesso}}{\#\text{tarefas elegíveis tentadas}},\\
 A_O&=\frac{\#\text{artefatos requeridos disponíveis}}{\#\text{oportunidades elegíveis}}.
\end{align}
Qualidade condicional ao artefato será apresentada junto de $A_O$ e TSR: ausência não receberá nota perfeita nem desaparecerá do denominador. Haverá TSR por fluxo, média micro e macro não ponderada. Sucesso longitudinal exigirá integridade das dependências críticas em todas as sessões; as cinco seeds serão primeiro agregadas por trajetória e depois por cluster com pesos definidos antes do teste.

Delegação e uso de ferramenta terão quatro níveis separados: parse/schema, seleção/intenção, argumentos/domínio e efeito observado. O sucesso de delegação será
\begin{equation}
 D=\frac{\#\text{artefatos válidos entregues e consumidos corretamente}}
 {\#\text{delegações elegíveis tentadas}}.
\end{equation}
Exigirá tipo/versão esperados, snapshot compatível e destinatário correto, com entrega e consumo no ledger. Apenas enviar mensagem não contará como sucesso; falhas de geração, schema ou receptor permanecerão no denominador. Também se medirão refs incompatíveis, reparos e terminais de handoff. Violação crítica será reportada por tentativa, artefato publicado e presença de qualquer violação por cluster, evitando diluição por muitas chamadas corretas.

Tokens serão totalizados em todas as gerações/reparos; latência incluirá fila, retrieval, geração, solver e validação. Espera humana será excluída do ativo e indicada separadamente. Serão apresentados p50/p95, taxa de timeout, conclusão censurada, prefill/decode quando disponíveis e RSS/VRAM. Quantil não identificável por censura será declarado, não estimado apenas pelos sucessos. Energia ficará N/A sem medidor; custo local terá segundos, tokens e amortização explícita, e custos de avaliação/juiz ficarão separados.

\subsection{Condições comparativas e justiça de recursos}

As condições serão MA completo, agente único equivalente, MA sem memória longitudinal, MA sem compactação e MA com todas as skills carregadas. O baseline receberá o mesmo Bonsai/runtime, corpus, ferramentas, memória disponível, validadores, solver, regras de persistência e capacidades; a diferença será um papel LLM em vez de cinco especializações. Seu prompt será ajustado nos mesmos dados e com esforço comparável, sem fragilidade deliberada. A ablação sem memória removerá apenas histórico longitudinal, preservando fontes e estado indispensável à rota corrente.

Um contraste focal adicional comparará \textbf{a mesma skill completa e reduzida}, distinguindo redução de instruções de seleção do catálogo. Objetivo, casos, modelo, fontes, ferramentas, schemas e limites serão os mesmos; versões dos dois prompts serão registradas. Serão propostos 60 casos distribuídos entre habilidades, duas variantes e cinco seeds (600 runs). Qualidade composta e custo total, incluindo reparos, serão pareados. Margem de preservação e família de hipóteses desse contraste serão justificadas e congeladas antes do teste; a margem de compactação não será automaticamente transferida para skills.

Todos terão janela 32.768, mesmas reservas, seis chamadas, oito ferramentas, 16 passos, reparos 2/1, 120 s ativos, geração de 90 s e solver de 10 s totais. Esses tetos serão ajustados apenas no piloto e depois congelados. Para distinguir especialização de aumento de inferência, curvas de qualidade--custo usarão níveis de chamadas 2/4/6 e entrada operacional de 8 mil/13.312 tokens, sem relaxar reservas. Serão reportados consumo real e tarefas impedidas pelo orçamento. Redução de custo sem preservação de qualidade não será descrita como eficiência bem-sucedida.

Seeds principais serão 0--4, e testes de falha com geração usarão 0--2; sem geração, teste determinístico. Gerador de dados, modelo, controlador e bootstrap terão RNGs separados. Cada execução restaurará snapshot idêntico em sandbox novo. Ordem será contrabalançada por bloco de caso/condição, com seed de shuffle registrada. Warmup e política de cache serão explícitos; não haverá cache de resposta gerada entre condições. Cold e warm serão separados, e o servidor local rodará sem concorrência não documentada. Offline será condição própria; conectado utilizará mocks/snapshots pareados e posterior smoke real do conector.

\subsection{Orçamento experimental e anotação}

O teste principal cruzará 120 consultas e 120 tentativas com cinco condições e cinco seeds, totalizando 6.000 runs. O longitudinal terá $30\times8\times5\times5=6\,000$ sessões, cada qual com vários turnos possíveis. O contraste de skills acrescentará 600 runs. Para falhas, dez casos por F01--F24, dois braços, duas condições (MA/SA) e três seeds darão até 2.880 runs, com menos repetições quando não houver geração. Agendas e smoke terão contagens separadas. Esses totais são planejamento, não execuções realizadas.

Antes do freeze, um piloto de 20 casos de desenvolvimento medirá wall time, chamadas e minutos de anotação. Para estratos $s$ com $N_s$ unidades previstas, segundos médios $\bar t_s$ e minutos médios de anotação $\bar a_s$, a projeção será
\begin{align}
 H_{exec}&=\frac{\sum_s N_s\bar t_s}{3\,600},\\
 H_{hum}&=\frac{2\sum_s N^{\mathrm{anot}}_s\bar a_s}{60}+H_{\text{adjudicação}}.
\end{align}
O fator 2 corresponde aos julgamentos independentes; $N^{\mathrm{anot}}_s$ seguirá a amostragem proposta, não o total de turnos automaticamente. Os autores aprovarão orçamento e ajustarão N com justificativa, ou classificarão análises como exploratórias, antes do teste final.

A proposta inicial de avaliação humana será 100\% dos candidatos críticos e amostra aleatória de 20\% dos restantes em cada estrato condição/disciplina/tipo. Probabilidade de inclusão, julgamentos e ICs amostrais serão preservados, com estimação ponderada quando necessária. Cobertura humana, disponibilidade de artefatos e número de clusters serão reportados. Gates semânticos se limitarão ao alcance dessa amostra e de sua inferência; não se afirmará verificação humana de 100\% das saídas. Checks determinísticos serão executados em todas as unidades elegíveis.

\subsection{Robustez, exposição e recuperação}

O \code{FaultController} injetará falhas por entrada, fronteira de ferramenta, processo ou ambiente, com alvo semântico, trigger, intensidade, duração e horizonte de recuperação. Guardará independentemente se o trigger foi alcançado e se a falha foi injetada. Teste de fronteira não será chamado de E2E: este exigirá o caminho real escolhido pelo agente. Exposição não alcançada será reportada e não contará como resistência demonstrada.

Cada caso perturbado terá par limpo com mesmo snapshot/seed, medindo $\Delta$qualidade, sucesso, custo, latência, efeitos indevidos e detecção/recuperação. Ataques terão taxa de sucesso por efeito (ASR), com denominadores elegível e exposto e utilidade legítima: bloquear tudo não será confundido com defesa útil. O desenho de objetivo legítimo/adversarial adapta o AgentDojo \cite{debenedetti2024agentdojo}; fluxos multiturno dialogam com AgentBench e BFCL \cite{liu2024agentbench,bfclMultiTurn}.

Crashpoints cobrirão antes/durante transação, depois de commit antes de checkpoint/ACK, após lease antes de envio, após efeito antes de ACK, após recibo e durante índice/migração/backup. \code{SIGKILL} testará crash de processo, não queda elétrica. Busy será provocado por conexão controlada; ENOSPC por volume isolado; ledger externo sobreviverá ao processo. Nenhuma falha será injetada no banco real de usuário. O horizonte inicial de recuperação será 300 s. Watchdog externo terá 120+5 s para encerrar processo; terminal após 120 s contará como overshoot, não novo SLO.

Concorrência será ensaiada com dois clientes na mesma revisão: um commit vence, outro recebe conflito; retry com ID/digest iguais recebe recibo antes da versão obsoleta; digest diferente é negado. Históricos de até oito operações serão checados por enumeração de ordens que respeitem invoke/return e modelo sequencial local. Checker de 5 s ou enumeração de 10 s esgotados produzirão inconclusivo. Linearizabilidade local não será estendida indevidamente aos efeitos remotos.

Testes stateful com Hypothesis gerarão sequências de tentativa, ajuda, avaliação, commit/retry, correção, aprovação, cancelamento, restart e delete/restore \cite{hypothesisStateful}. Invariantes serão conferidos após cada ação, inclusive chamadas proibidas. Metamorfismos determinísticos incluirão permutação de IDs/demandas preservando viabilidade/ótimo, mais disponibilidade não removendo solução, mais bloqueios não criando solução e repetição de operação preservando efeito \cite{chen2018metamorphic}. Paráfrase dependerá de ouro humano; ruído e distrator serão sensibilidade, não invariância de ranking garantida.

Falhas compostas prioritárias serão OCR--data--agenda; falso crédito--redutor--compactação; mudança de snapshot--crash; correção/exclusão--resumo antigo; loop/replay--efeito; e timeout--trace--viés de análise. No desenho limpo/A/B/A+B, a interação será $Y_{AB}-Y_A-Y_B+Y_{limpo}$, com exposição de ambas as causas conferida. A seleção será por severidade, sem alegar cobertura de todas as combinações.

\subsection{Julgamento humano, inferência e decisões}

Artefatos semânticos serão exportados cegos, com pergunta, fontes e rubrica, sem condição, trace de identidade ou parecer do Monitor. Dois avaliadores julgarão independentemente nas unidades selecionadas; acordo será medido antes da adjudicação. Juiz LLM será secundário e calibrado em validação, com modelo/prompt/hash separados. Comparações inverterão ordem A/B, controlarão comprimento sem conteúdo novo e paráfrases, registrando flips e erros críticos \cite{zheng2023judge}. Modelo diferente não garante independência do ouro nem amplia o alcance de um gate humano amostral.

As primárias serão TSR longitudinal e violação crítica publicada, com decisão conjunta. A regra proposta para benefício será limite inferior do IC95\% da diferença pareada de TSR (MA menos agente único) acima de zero e nenhuma violação crítica publicada no MA; a comparação de violações entre condições e sua incerteza será reportada conjuntamente. Não haverá tolerância de engenharia para violação crítica; observação de violação bloqueará a capacidade afetada, e ausência em amostra finita não provará risco nulo nem equivalência de riscos. Família de secundárias confirmatórias, incluindo sustentação, memória e custo, terá correção de Holm; análises adicionais serão exploratórias. Estimandos micro/macro, pesos e hipóteses serão definidos antes do teste, alinhando análise ao desenho \cite{dror2018hitchhiker}.

Será proposto bootstrap pareado por cluster, inicialmente \textbf{10.000 reamostragens}, seed 2026, mantendo conjuntamente casos, seeds, condições e braços. A cada reamostragem serão recalculadas razões micro ou médias macro conforme o estimando. IC95\% percentil será default a validar por simulação no piloto, não BCa aplicado a turnos independentes. O agrupamento exige implementação explícita; \code{paired=True} sozinho não resolve clusters \cite{field2007cluster,scipyBootstrap}. Poucos componentes e zero eventos poderão produzir IC degenerado; esse limite será declarado.

Para compactação, a margem inicial proposta de não inferioridade será \textbf{0,03} na taxa composta de emitir artefato e passar todos os essenciais. É tolerância de engenharia a justificar com especialistas antes do congelamento, não margem extraída da literatura. Com $\Delta=Q_{\text{compactação}}-Q_{\text{sem compactação}}$, preservação exigirá limite inferior unilateral de 95\% acima de $-0{,}03$, custo menor e nenhuma tolerância em violação crítica. Valor $p>0{,}05$ não demonstrará equivalência. Tanto margem quanto 10.000 reamostragens serão propostas verificadas antes do freeze, sem ajuste orientado pelo teste final.

Piloto apenas em desenvolvimento simulará diferenças pareadas por cluster em grades de variância, correlação e número de componentes, estimando potência e precisão para efeito relevante definido com os responsáveis \cite{lakens2022inferences}. Se 30 trajetórias forem insuficientes, ampliar-se-á a coleta antes do congelamento ou a análise será declarada exploratória por recursos. Não haverá potência pós-hoc. Com zero eventos e $G$ clusters independentes comparáveis amostrados, o limite unilateral de 95\% para presença de qualquer violação será $1-0{,}05^{1/G}$; não será aplicado a ataques adversariais escolhidos nem a turnos correlacionados.

\subsection{Execução, validade e produtos}

O pré-flight conferirá autorização/splits, hashes, admissão de modelo/ambiente, controles dos oráculos, smoke offline e restore. O runner restaurará sandbox, iniciará ledger e relógios, executará roteiro até terminal/budget, capturará domínio/efeitos/traces e exportará artefatos cegos. Relógio virtual controlará datas e lógica; relógio monotônico real conferirá backend/deadlines. Após anotação serão calculados acordo, métricas, clusters, ICs e decisões. Freeze produzirá manifesto imutável com dados, ouro, scripts/ataques, configs, seeds, thresholds, margens, budgets, hipóteses, parsers e oráculos. Reprodutibilidade seguirá inventário de pesos/binários e ensaio offline, não apenas um lock de pacotes \cite{pineau2021reproducibility}.

O relatório experimental informará N planejado, executado, elegível e exposto, cobertura W/O/F, qualidade e disponibilidade, TSR, violações, utilidade/ASR, clean/fault, recuperação/censura, custo, ablações, ICs e acordo/meta-avaliação. Falha crítica bloqueará a liberação da capacidade afetada, preservando todas as runs. Threshold pontual com IC amplo será evidência limitada. Esta fase descreve metodologia e passos de implementação; não contém testes executados nem resultados de eficácia.

A validade interna depende de snapshots, ouro independente, orçamento equivalente e cegamento; a validade de construto depende de não confundir fluência, acerto assistido, domínio e contenção segura com sucesso original. Componentes independentes do harness reduzem, mas não eliminam, erros correlacionados. Duas disciplinas, trajetórias roteirizadas, corpus finito e um hardware limitam validade externa. Simulações avaliam consistência do sistema, não ganho humano. Estudo educacional futuro precisaria de participantes, pré/pós-teste tardio, transferência, controle de tempo e protocolo próprio. O pacote de avaliação será concebido para tornar falhas e limitações observáveis, em vez de preencher lacunas com resultados presumidos.

\section*{Apêndice}
\label{sec:apendice}

\subsection*{A. Rastreabilidade das famílias de fluxo}

Os IDs abaixo conectam contratos, casos e relatório experimental. ``Cobertura'' significará presença de casos executados e elegíveis, não apenas existência de uma linha no catálogo.

\begin{center}\small
\begin{longtable}{@{}L{1.0cm}L{5.7cm}L{2.8cm}L{5.3cm}@{}}
\toprule ID & Família & Saídas & Terminal/variante essencial \\ \midrule
\endfirsthead
\toprule ID & Família & Saídas & Terminal/variante essencial \\ \midrule
\endhead
W01 & Ingestão e confirmação & O01--02 & Indexado, pendência, ilegível, revisão \\
W02 & Consulta fundamentada & O03--04 & Apoio, parcial, geral rotulada, abstenção \\
W03 & Prática e apoio gradual & O05--07 & Espera, pistas, solução, recusa \\
W04 & Avaliação formativa & O08--09 & Aceita/abstenção, ajuda preservada \\
W05 & Retomada longitudinal & O10--11 & Atual/histórico, correção, ausência \\
W06 & Planejamento & O12--13 & Viável, conflito, inconclusivo, ambíguo \\
W07 & Replanejamento & O09, O13 & Diff, fixos, déficit, versão nova \\
W08 & Exportação/sincronização & O14 & ICS, recibo, efeito desconhecido \\
W09 & Curadoria opcional & O15 & Texto, transcrição, metadados ou offline \\
W10 & Consulta/correção/exclusão & O10, O16 & Confirmação, propagação, conflito \\
W11 & Recuperação operacional & O16--17 & Replay, cancelado, limite, reconciliação \\
W12 & Instalação e manutenção & O18 & Doctor, backup/restore, smoke offline \\
\bottomrule
\end{longtable}
\end{center}

\subsection*{B. Matriz completa das 18 saídas}

Os limites desta matriz são \textbf{propostos} e serão congelados após desenvolvimento/validação. Percentuais de 100\% e zero violações são critérios no ensaio finito; estimativas, disponibilidade, denominadores e ICs serão reportados conjuntamente.

\begin{center}\small
\begin{longtable}{@{}L{0.8cm}L{3.0cm}L{6.3cm}L{4.7cm}@{}}
\toprule ID & Artefato & Oráculo independente/métrica & Gate inicial proposto \\ \midrule
\endfirsthead
\toprule ID & Artefato & Oráculo independente/métrica & Gate inicial proposto \\ \midrule
\endhead
O01 & Manifesto e blocos & Transcrição humana; CER/WER, exact match crítico, pares de ordem e resolução de página/bbox/versão & EM crítico $\geq0{,}98$; refs corretas 100\% \\
O02 & Tópicos, fatos e prazos & Entidades/relações 1:1; tupla de data, tipo, precisão, zona, autoridade e status; ledger de confirmação & $F_1\geq0{,}90$; promoção ambígua silenciosa zero \\
O03 & Evidências e citações & Qrels original--chunk; recall@40, MRR/nDCG@6, aspectos, CP/CR; negativos de suporte & Recall $\geq0{,}90$; CP $\geq0{,}95$; abstenção separada \\
O04 & Explicação & Humanos: corretude, pertinência, clareza e limites 0--4; claims necessárias e erro crítico & Mediana $\geq3$ por dimensão; crítico bloqueia caso \\
O05 & Exercício/quiz & Especialista resolve antes do gabarito; tolerância/unidade e equivalência simbólica; fonte/solução separadas & Validade $\geq0{,}95$; gabarito sem origem rejeitado \\
O06 & Flashcard & Checklist de atomicidade, precisão, frente/verso e fonte; duplicata por objetivo e reveal & Validade $\geq0{,}95$; revelação antecipada zero \\
O07 & Turno/pista & Rubrica de contingência/participação 0--4 e máquina de estados externa; chamadas durante espera & Mediana $\geq3$; terminais respeitados 100\% \\
O08 & Avaliação e evidência & Critérios por tópico 0--2; present/missing, kappa, macro-$F_1$, falso crédito/aceitas e críticos & $\kappa\geq0{,}70$; falso crédito $\leq0{,}05$; cobertura amostral \\
O09 & Estado/revisão & Replay por julgamento vigente, ordinal e famílias; revisão tardia, dedup e independência por sessão & Replay 100\%; favorável fiel; estimativa nula \\
O10 & Memória temporal & Ledger bitemporal e QA atual, retrospectiva, multissessão e abstenção; cortes de conhecimento & Acurácia $\geq0{,}90$; vazamento zero; correções respeitadas \\
O11 & Síntese/contexto & Tuplas críticas, claims--eventos e tokenizer após template; contradições e qualidade downstream & Prazos/ajuda críticos 100\%; budget válido; fonte exigida \\
O12 & Demanda de estudo & Ouro de restrições originais; omissões/duras extras, status da duração e pertinência & Duras fiéis 100\%; nenhuma inventada; estimativa rotulada \\
O13 & Plano, conflito e diff & Checker em minutos UTC originais; enumeração pequena, vetor objetivo e diff por bloco & Violação publicada zero; status/diff corretos; grade explicitada \\
O14 & ICS/recibo remoto & Parser distinto, UTF-8/CRLF/folding, UID/SEQUENCE/timestamps e tuplas plano--ICS--readback & Equivalência 100\%; duplicata zero no adaptador ensaiado \\
O15 & Recurso/estado de acesso & Snapshot de fetch/transcrição, proxy de acesso e pertinência humana & Pertinência $\geq0{,}90$; resumo não lido zero \\
O16 & Recibo, alteração e exclusão & Ledger, revisão global/aprovação, canários, dependentes e restore com tombstones & Duplicata/perda zero; exclusão lógica e retenção conferidas \\
O17 & Trace/medidas & Proxy externo versus spans; chamadas/reparos/tokens, tempos e recursos; censura & Spans obrigatórios 100\%; falhas incluídas; ausente nulo \\
O18 & Bundle e diagnóstico & Inventário e hashes; instalação offline limpa, smoke, integridade/FK e restore lógico & Artefatos requeridos presentes; smoke/restore íntegros \\
\bottomrule
\end{longtable}
\end{center}

\subsection*{C. Matriz completa das 24 falhas}

Cada linha será instanciada por plano de falha com trigger e exposição conferidos. A medida de contenção será apresentada ao lado da utilidade original e dos deltas limpo/perturbado.

\begin{center}\small
\begin{longtable}{@{}L{0.8cm}L{3.0cm}L{6.3cm}L{4.7cm}@{}}
\toprule ID & Fronteira & Injeção concreta & Observação/oráculo \\ \midrule
\endfirsthead
\toprule ID & Fronteira & Injeção concreta & Observação/oráculo \\ \midrule
\endhead
F01 & Extração/OCR & Scan com rotação/ruído; ordem, sinal, número/data mutados e anotados & O01--02, gabarito/plano afetados; diagnóstico versus promoção \\
F02 & Ontologia/data & Prazo parcial/relativo, revisão conflituosa e alias ambíguo & O02/O12; falsa confirmação e restrição downstream \\
F03 & Recuperação & Remover fonte, acrescentar distrator e apresentar versão antiga & O03; recall, apoio, abstenção e disponibilidade \\
F04 & Geração/citação & Claim sem apoio, fonte contraditória, citação existente irrelevante & O03--04; falso suporte e erro factual \\
F05 & Tutoria/ajuda & Scripts de solução/recusa, três pistas, reveal e resposta posterior & O05--07/O09; vazamento, insistência, independência falsa \\
F06 & Monitor/gabarito & Curta correta, longa errada, paráfrase e referência defeituosa & O08--09; falso crédito, risco--cobertura, estado \\
F07 & Redutor & Retry de tentativa, ajuda, replay com futuro e novo erro crítico & O09; duplicata, corte temporal e falso favorável \\
F08 & Memória & Correção retroativa, homônimo, supersedes e evidência antiga & O10; atual/histórico, origem e conflito \\
F09 & Compactação & Pressão de contexto e síntese com data/polaridade/ajuda alterada & O11 e downstream; retenção crítica e rejeição \\
F10 & Comunicação & Envelope/schema/ref/usuário incompatíveis, inclusive JSON válido & Domínio e validadores; aceitação inválida, reparos/efeitos \\
F11 & Grafo/tools & Erro persistente de ferramenta, rota circular e pedidos repetidos & Watchdog; calls/steps, overshoot e terminal \\
F12 & Inferência e recursos & Timeout, truncamento, OOM, tokenizer divergente e fila saturada & Disponibilidade, censura, erro e promoção parcial \\
F13 & Solver/tempo & 17 min reduzidos a 15, meio-slot, DST, ciclo e UNKNOWN & O13; checker original, status e publicação inválida \\
F14 & Versão e replanejamento & Alterar snapshot antes de commit, dois clientes e bloco iniciado & O13/O16; conflito, lost update, churn/fixo \\
F15 & Efeito/retry & Crash nas barreiras de commit, lease, efeito, ACK/recibo; timeout remoto & Ledger externo; duplicata/perda, reconciliação, unknown effect \\
F16 & Prompt injection & Instrução em PDF/site/tool buscando mudar prazo ou executar efeito & Proxy e domínio; ASR por efeito e utilidade legítima \\
F17 & Isolamento e acesso & Outro usuário, traversal/symlink, DNS/redirect privado e canário & Leituras/envios/efeitos negados; vazamento zero \\
F18 & Persistência & Busy, ENOSPC, crash em migração, backup ou índice incompleto & O16/O18; atomicidade antigo/novo, restore e recuperação \\
F19 & Web/conector & Offline, 404, quota, mudança de recurso e falta de transcrição & O15; continuidade local e acesso fiel \\
F20 & ICS e interoperabilidade & UTF-8, DATE/all-day, zona, UID e revisão/DTSTAMP inconsistentes & O14; parser distinto, equivalência e retry \\
F21 & Correção e exclusão & Excluir/corrigir e depois reindexar, replay e restaurar backup & O10/O16; dado antigo/canário não ressurgem \\
F22 & Telemetria & Span omitido, retry invisível, contador/exporter ausente & O17; proxy externo e erro de medição, nulo explícito \\
F23 & Validade experimental & Ordem/comprimento do juiz e caso deliberadamente vazado & Auditor/controle humano; análise contaminada sinalizada \\
F24 & Empacotamento & Remover peso, OCR/lock; hash errado; instalação offline limpa & O18; doctor, dependência ausente e smoke/restore \\
\bottomrule
\end{longtable}
\end{center}

\subsection*{D. Declaração de uso de IA generativa}

A elaboração deste relatório foi assistida pelo OpenCode na leitura e consolidação dos documentos do projeto, na redação técnica e na formatação LaTeX. A ferramenta foi utilizada como apoio à elaboração. A conferência final para submissão e a responsabilidade pelo conteúdo cabem aos autores identificados na capa.

\newpage
\IfFileExists{aasjournal.bst}{\bibliographystyle{aasjournal}}{\bibliographystyle{plainnat}}
\bibliography{references}

\end{document}
